% =============================================================================
% Title:        Gas tubing system
% Author:       Nathaniel Thomas
% Affiliation:  Texas A&M University, Department of Nuclear Engineering
% Email:        nathaniel@tamu.edu
% Date Created: 3/16/2026
% Version:      1.0
%
% Description:
%   A .tex source file describing the tubing system of the purifier.
%
% =============================================================================

\section{Tubing system}

This section describes the tubing system built to introduce gases into the
FLiNaK purifier. In addition, it also describes the scrubber system built to
safely capture and dispose of any residual \ch{HF} that was not
reacted during the sparging process.
A diagram of the gas mixing system is shown in \Cref{fig: gas mixing
diagram}. \par

% TODO: Find a diagram of this system
\begin{figure}
  \centering
  % \includegraphics[width=0.6\textwidth]{}
  \caption{A diagram of the gas mixing system for the FLiNaK purifier.}
  \label{fig: gas mixing diagram}
\end{figure}

\subsection{Gas sparging system}

The gas sparging system is meant to introduce gas into the
molten salt, resulting in the chemical reactions to enable
purification, as described in \Cref{sec: hf sparging}.
There are three gases used in total:

\begin{enumerate}
  \item Hydrogen gas
  \item Hydrofluoric acid
  \item Argon gas
\end{enumerate}

Each of the gases provide a different purpose for purification.
Argon removes residual water, hydrofluoric acid converts metallic
oxides to water and metallic fluorides, and  hydrogen reduces
metallic impurities in the salt into their elemental forms, causing
them to fall out of solution and be filtered out during the transfer
process. An image of the gas feed system is shown in
\Cref{fig: gas control system,fig: gas control inside}.
\footnote{Note that the filtration step has been removed in
  the current iteration of the design in order to mitigate issues
  during the transfer process. These elemental metallic impurities
will have to be filtered out separately.} \par

% TODO: Find or make diagram of this system
\begin{figure}
  \centering
  % =============================================================================
% Title:        Gas control system
% Author:       Nathaniel Thomas
% Affiliation:  Texas A&M University, Department of Nuclear Engineering
% Email:        nathaniel@tamu.edu
% Date Created: 3/18/2026
% Version:      1.0
%
% Description:
%   A .tex source file showing an annotated image of the FLiNaK purifier gas
%   control tubing outside the purifier.
%
% =============================================================================

\begin{tikzpicture}

  % Place image as a node
  \node[anchor=south west, inner sep=0] (img) at (0,0)
  {
    \includegraphics[width=10cm]{/home/nathaniel/Code/LaTeX/Research/flinak-purifier-document/assets/pictures/gas
    control system.jpg}
  };

  % Define coordinate system relative to image
  \begin{scope}[x={(img.south east)}, y={(img.north west)}]

    % Example annotations (coordinates are now 0–1 normalized)
    \draw[->, red, thick] (0.5,0.7) -- (0.63,0.8)
    node[red,pos=-0.3,above] {Moisture filter};

    \draw[->, red, thick] (0.5,0.5) -- (0.63,0.5)
    node[red,pos=-0.3,above] {Ar flow controller};

    \draw[->, red, thick] (0.6,0.35) -- (0.6,0.23)
    node[red,pos=-0.1,above] {Ar flow controller};

    \draw[->, red, thick] (0.3,0.25) -- (0.43,0.25)
    node[red,pos=-0.3,above] {Ar purge valve};

  \end{scope}

\end{tikzpicture}

  \caption{An annotated image of the gas mixing system outside of the
  fume hood (11/22/2022)}
  \label{fig: gas control system}
\end{figure}

\begin{figure}
  \centering
  \input{./src/d&c/tubing/gas_control_inside.tex}
  \caption{An annotated image of the gas mixing system inside of the
  fume hood (11/22/2022)}
  \label{fig: gas control inside}
\end{figure}

\subsection{Scrubber system}

This system is designed to safely dispose of the residual gases from the
gas sparging system. The scrubber gas is collected from the reaction chamber
following purification. There is a dedicated port on the top of the reaction
chamber. This port then feeds into a tube that runs to a solid scrubber
containing calcium oxide. This neutralizes some of the hydrofluoric acid,
reducing the load on the wet scrubber.
\footnote{This tube has a tee that feeds
into the RGA system, enabling analysis of the scrubber gas content.}
Following the dry scrubber, the line feeds into a KOH wet scrubber, where
the remaining HF is neutralized.
\par

% TODO: Include an image of the scrubber solution.
